\scanpage{155}
% archive sheet number: 154
% L1–2
On aimerait scinder cette extension. \deleted{Notons}\inserted{Posons}
\[
\mathrm{Et}_F(s_0)\overset{\mathrm{def}}{=}\coprod_{i\in\Pi}\mathrm{Et}_i(s_0)
=\{s\in S\mid\{s_0,s\}\in A_0\}\,,
\tag*{(72)}
\]
% L3–4
et notons
\[
\deleted{\unsure{G\ \text{est}}}\ G_{s_0}\text{ est transitif sur }\mathrm{Et}_F(s_0)
\tag*{(73)}
\]
% L5–10
ce qui résulte encore formellement du fait que $G$ est transitif sur $\vec A$, \deleted{et $G_0^+$ est transitif sur} le stabilisateur \inserted{dans $G_{s_0}$} d'un élément $t_0\in\mathrm{Et}_F(s_0)$, i.e. le stabilisateur dans $G$ de $a\ (s_0,t_0)\in\vec A$, est un sous-groupe d'ordre 2, contenu dans $G_0^+$ \deleted{\unsure{donc dans}}, de ``déterminant'' égal à $-1$. Son élément non trivial mérite la notation $\tau_\infty$, plus précisément $\tau_\infty^{a_0}$ où $a_0=\{s_0,t_0\}$, égal à $\tau_\infty^{r_0}$ où $r_0$ est un repère de la forme
% L14
\[
r_0=(s_0,t_0,u_0)\ .
\]

% L15–17
Notons que le groupe
\[
Z\simeq(\mathbb Z/n\mathbb Z)^*
\]
opère \uline{librement} sur $S$ (puisqu'il opère par autom d'une structure d'espace homogène), soit $D_t$ \deleted{\unsure{soit}} l'orbite \deleted{\unsure{sous}} d'un $t\in S$. Notons que $Z$ stabilise $\mathrm{Et}_F(s_0)$ (pour tt $s_0\in S$), donc toute orbite \inserted{de $Z$ (un ``\uline{rayon}'')} qui rencontre $\mathrm{Et}_F(s_0)$, est contenue dedans. Choisissons donc, en plus de $s_0\in S$, un rayon \deleted{\unsure{$D$}}
% L23–24
\[
D_0\deleted{\unsure{=Z.t_0}}\subset\mathrm{Et}_F(s_0)\ ,\qquad D_0\ \inserted{Z\text{-}}\text{orbite d'un }
t_0\in\mathrm{Et}_F(s_0)
\tag*{(74)}
\]
